\section{Agent School 与可靠性工程：让系统持续变强}

Virtual Lab 的长期价值依赖于``能否自我更新''。课程中提到的 Agent School 可以理解为研究团队的持续培训机制：指定主题、检索文献、学习工具、回到任务中应用。其重点不是一次训练，而是\textbf{能力增量可累积}。

\subsection{Agent School 的训练单元}

从字幕内容看，训练主题包括 AlphaFold 等新工具使用、跨领域知识补齐、任务相关文献快速吸收。人类提供关键主题和边界，Agent 执行学习与整合。这是典型的人机分工：人定义方向，机器承担高吞吐知识摄取与初步结构化。

\begin{knowledgebox}{Agent School 的输入/输出接口}
\begin{itemize}
    \item 输入：学习主题、约束条件、目标任务上下文。
    \item 中间产物：文献摘要、工具使用脚本、失败案例归档。
    \item 输出：可被 PI 调度的新角色能力或升级版工作流。
\end{itemize}
\end{knowledgebox}

\subsection{外部记忆与 sandbox：从会话体到工程体}

Zou 在问答中多次强调外部 memory 与 sandbox。memory 让 Agent 不必反复丢失上下文，sandbox 让工具调用可控可复现实验。二者结合后，系统从``聊天模型''升级为``研究执行体''。

\begin{importantbox}{memory + sandbox 的组合价值}
\begin{itemize}
    \item memory 负责长期一致性：保存假设、证据、反例与版本。
    \item sandbox 负责执行可靠性：固定依赖、记录命令、可重放运行。
    \item 二者共同降低``看起来懂''但``做不出来''的风险。
\end{itemize}
\end{importantbox}

图\ref{fig:vl-school} 展示了 Agent School 的关键页。\footnote{来源：Berkeley RDI 课程视频，关键帧时间戳 00:14:40。}

\begin{figure}[H]
\centering
\includegraphics[width=0.92\textwidth]{frame-01-virtual-lab.jpg}
\caption{Agent School：选题、生成查询、检索与筛选论文}
\label{fig:vl-school}
\end{figure}

\begin{knowledgebox}{读图：Agent School 是能力更新流水线}
左侧先选专家人格与学习主题，中间由问题生成器展开知识缺口，右侧检索并筛选论文。输出不应只是摘要，而应进入可测试的工具脚本、失败案例和角色能力版本。\textbf{老师强调：}长期记忆与 sandbox 让这些更新可复用、可回滚，否则“学习”只停留在一次会话。
\end{knowledgebox}

\subsection{Human-in-the-loop 的真实作用}

课程给出非常务实的表述：人类不必逐句监督 Agent，但要在关键节点干预，防止 topic drift 和误设前提。尤其在开放性科学任务中，``问题定义本身''就是高价值决策，不能完全外包。

\begin{warningbox}{把人类放在末端验收会导致系统性风险}
如果人类只在最终结果阶段介入，前期错误假设可能已经驱动大量无效计算与错误实验优先级。更稳妥做法是在选题、假设变更、关键工具切换、外部验证前后设置人工闸门。
\end{warningbox}

\subsection{人格多样性与并行会议}

Zou 还提到给不同 Agent 配置不同人格和偏好，在并行会议里比较结果。其本质是制造有控制的``认知异质性''，避免单一推理路径导致的脆弱性。对于科研而言，这等价于把``不同导师风格/学科传统''程序化复用。

\begin{importantbox}{可操作实践：并行人格实验}
\begin{itemize}
    \item 为同一课题设置保守、激进、成本敏感三种 PI 策略。
    \item 对比三组输出在新颖性、可执行性、验证成本上的差异。
    \item 让 Critic 对每组给出失败模式预测，再做合并决策。
\end{itemize}
\end{importantbox}

\subsection{本章小结}

Agent School、memory、sandbox、HITL、人格多样性一起构成了系统长期可靠性的骨架。它们解释了为什么某些多 Agent 系统``越跑越稳''，另一些系统``越跑越飘''。

